%=============================================================================!
% 16.3  BONDS AND MOLECULES                                                   !
%=============================================================================!
\section{Bonds and molecules}
\label{sec:ob-bonds}

To follow a reaction, we must first decide which atoms form each
molecule. We assign bonds to pairs, then collect atoms joined through
those bonds into molecules. A distance cutoff gives a first rule, but
thermal vibrations can carry a pair across it repeatedly without a
chemical change. This section introduces a bond criterion that retains
that distinction, then uses the resulting graph to identify molecules
and species.

\subsection*{A bond from a distance}

The forces of most models say nothing about bonds: they are functions of
the positions. A bond is therefore assigned from the distance between
two atoms, and a common yardstick is the sum of their covalent radii,
$R_i + R_j$, from a table of radii for each element such as the one ASE
carries\cite{Larsen2017}: two carbon atoms ($R = \SI{0.76}{\angstrom}$
each) are bonded within some multiple of \SI{1.52}{\angstrom}. Read
literally, such a rule makes a bond flicker. A bonded pair vibrates, and
if the far end of its vibration lies near the cutoff, each vibration
carries it out and back across it, two events per period with nothing
chemical happening.

The cure is a band of hysteresis, as in a thermostat that switches the
heating on below one temperature and off above a higher one. A pair not
yet bonded becomes bonded when it comes closer than $r_{\mathrm{on}}$; a
bonded pair stays bonded until it parts beyond $r_{\mathrm{off}} >
r_{\mathrm{on}}$. A vibration whose whole swing, from its near end to its
far end, is narrower than the band cannot carry the pair across both
distances, so it cannot make a bond flicker, while a real dissociation,
which carries the atoms apart for good, crosses both. \texttt{mdlab.bonds.BondTracker} takes
$r_{\mathrm{on}}$ and $r_{\mathrm{off}}$ as 1.2 and 1.4 times the sum of
the covalent radii unless told otherwise, factors that bracket the first
minimum of $g(r)$ for carbon (\cref{sec:ob-reactions}); for a model, or
for any pair of elements, the two distances can be given directly.

To measure the effect we use a model gas built to react. Two kinds of
atom, A and B, of nitrogen's mass, attract each other through the Morse
well of \cref{sec:pe-pairs} with depth $D_{\mathrm M} = \SI{1}{\electronvolt}$,
width $a_{\mathrm M} = \SI{2}{\per\angstrom}$ and $r_0 = \SI{1.2}{\angstrom}$, the $D$
and $a$ of that section marked here to keep them apart from the
diffusion coefficient and the lattice spacing of later sections; two atoms of
the same kind only repel, as $D_{\mathrm M}\mathrm{e}^{-2a_{\mathrm M}(r - 2.6\,\text{\AA})}$, which
weakens a second partner's hold: a line B-A-B with both bonds at
\SI{1.48}{\angstrom} lies only \SI{0.40}{\electronvolt} below AB and a free
B, so such complexes form but soon part. The constants are chosen for
illustration and are not fitted to any real pair of atoms. 250 molecules
AB in a cubic cell of side \SI{63}{\angstrom} are held at 1500, 1750,
2000, 2250 and \SI{2500}{\kelvin} by Langevin friction of \SI{1}{\per\pico\second}, which
stands for a buffer gas that supplies and removes energy, for
\SI{100}{\pico\second} each. A bond forms where the pair's energy falls
below $-D_{\mathrm M}/2$, at $r_{\mathrm{on}} = \SI{1.814}{\angstrom}$, and breaks
where it rises above $-D_{\mathrm M}/8$, at $r_{\mathrm{off}} = \SI{2.570}{\angstrom}$.
A molecule's small vibrations have the period $2\pi\sqrt{m_{\mathrm r}/2D_{\mathrm M}a_{\mathrm M}^2}
= \SI{59.8}{\femto\second}$ (\cref{sec:os-small,sec:os-bodies}), with
$m_{\mathrm r} = \SI{7.0035}{\amu}$ the reduced mass and $2D_{\mathrm M}a_{\mathrm M}^2 =
\SI{8}{\electronvolt\per\angstrom\squared}$ the stiffness of the well.

At \SI{2000}{\kelvin} a single cutoff at $r_{\mathrm{on}}$ records 53\,265
bonds formed or broken in \SI{100}{\pico\second}; the band records 16\,442,
3.2 times fewer (\cref{fig:ob-bonds}a). Many of the remaining events are
not flicker either: an A and a B that meet with too much energy to stay
together pass through each other's well and part again within a
vibration or two. Requiring that a bond last a set time before its
forming and breaking count, and likewise that a pair which parts stay
apart that long, and dropping both events of any bond or gap shorter than
that (\texttt{bonds.persistent}), removes these passing encounters and the
flicker the band leaves. For a least lifetime of \SI{0.5}{\pico\second}
the counts fall to \num{5501} with the single cutoff and \num{6548} with
the band, which then keeps a fifth more. The remaining events fall
more slowly as the least lifetime grows, by 40\% and 36\% from 0.5 to
\SI{2}{\pico\second}, against a fall by 9.7 and 2.5 times from no least
lifetime to \SI{0.5}{\pico\second}: they are the molecules that form and
break.

\subsection*{Keys and set differences}

Each bond is stored as one whole number, its key $iN + j$ for the
atoms $i < j$. Different pairs have different keys, since $j < N$, and
the atoms are recovered as the quotient and the remainder of the key
divided by $N$. A computer stores such numbers in 64 binary digits, bits,
reaching $2^{63} \approx \num{9.2e18}$ with one bit for the sign, so keys
fit for $N$ up to about $3\times10^9$, and about $2\times10^9$ once the
tracker doubles each key (below). The
bonds formed between two frames are the keys in the new set and not in
the old, and those broken the keys in the old and not in the new.

\begin{toolbox}[Set differences of sorted lists]
Two lists of whole numbers, each sorted and without repeats, can be
compared in one pass. A pointer starts at the head of each. If the two
numbers under the pointers are equal, the number is in both lists, and
both pointers move on; otherwise the smaller is in its own list only,
and only its pointer moves. When one list is used up, the rest of the
other is in that list only. Each step moves at least one pointer, so for
lists of lengths $n_1$ and $n_2$ the pass takes at most $n_1 + n_2$
steps, against $n_1n_2$ for testing every pair. \texttt{BondTracker.update} merges this way
the sorted keys within $r_{\mathrm{off}}$ of the new frame with the
sorted bonds of the last one: a key in both is kept, a key only in the new
list is a new bond if it lies within $r_{\mathrm{on}}$, and a key only in
the old list is a bond broken.
\end{toolbox}

The candidates are the pairs that the force evaluation has already
found, which the model of \texttt{md.PairModel} keeps with their
distances after each call, so the tracker measures no distance of its
own; this serves as long as the cutoff of the forces reaches the largest
$r_{\mathrm{off}}$. Most candidates lie beyond the largest
$r_{\mathrm{off}}$ and are passed over at once; the rest are sorted, each
stored as twice its key plus 1 if the pair is within $r_{\mathrm{on}}$, so
that one sort orders the keys, and halving, the remainder dropped,
recovers each key, the remainder its flag; then they are merged:

\begin{code}
    def update(
        self, i: ArrayLike, j: ArrayLike, dist: ArrayLike
    ) -> tuple[NDArray, NDArray]:
        # the keys within r_off, each carrying its "closer than r_on" flag
        # in an extra lowest bit, so that one sort orders both
        packed = np.sort(_jit(_held_keys)(
            np.asarray(i, dtype=np.int64), np.asarray(j, dtype=np.int64),
            np.asarray(dist, dtype=float), self.kind, self.r_on, self.r_off,
            self._reach, self.n))
        self.keys, formed, broken = _jit(_merge)(
            packed >> 1, (packed & 1).astype(np.bool_), self.keys)
        return formed, broken
\end{code}

\noindent The two helpers are plain loops, the filter and the merge of the
toolbox, compiled by Numba as in \cref{sec:md-speed}.

\subsection*{Molecules}

A molecule is a set of atoms joined to one another through chains of
bonds.

\begin{toolbox}[Graphs and their connected components]
A graph is a set of nodes and a set of edges, each edge joining two
nodes. Here the atoms are the nodes and the bonds the edges. A graph of
$N$ nodes can be stored as its adjacency matrix, $N\times N$, with 1 in
row $i$ and column $j$ when $i$ and $j$ are joined and 0 otherwise. For
atoms nearly every entry is 0, since each atom has a few bonds, so only
the positions of the ones are stored, a sparse matrix whose storage grows
as the number of bonds, not as $N^2$. A connected component is a largest
set of nodes any two of which are joined by a path of edges. It is found
by a search, which starts from any node not yet assigned, adds its
neighbours, then their neighbours, and so on until no new node is
reached; those nodes are one component, and the search starts again from
a node not yet assigned. Each edge is followed at most twice, so the work grows
as the number of nodes plus edges.
\end{toolbox}

\noindent \texttt{mdlab.bonds.molecules} builds the sparse matrix of the
bonds and labels each atom with its component through
\texttt{scipy.sparse.csgraph}:

\begin{code}
def molecules(n_atoms: int, keys: ArrayLike) -> NDArray:
    i, j = np.divmod(np.asarray(keys, dtype=np.int64), n_atoms)
    graph = coo_matrix((np.ones(len(i)), (i, j)), shape=(n_atoms, n_atoms))
    return connected_components(graph, directed=False)[1]
\end{code}

\subsection*{Species}

Two molecules belong to the same species if they have the same atoms
bonded in the same way, whatever the numbers the atoms carry in the run.
The formula is the first part of the name: the number of atoms of each
element, written in the Hill order: in a molecule with carbon, carbon
first, then hydrogen, then the others alphabetically, and in one without,
all of them alphabetically, so that a molecule of two A and one B is
A$_2$B.
The formula does not tell isomers apart, molecules of the same atoms
bonded differently, such as a ring of six carbon atoms and a chain of
six.

\begin{toolbox}[Hashing and canonical labels]
A hash is a function that turns data of any size into a number of fixed
size, such that equal data give equal numbers and different data give
different numbers with high probability: a fingerprint. Two molecules
should get the same number if one can be renumbered into the other. A
label that does not change under renumbering is called canonical; sums
over the atoms are canonical, since the order of the terms does not
change a sum. \texttt{mdlab.bonds.graph\_hash} builds one in rounds, after
Weisfeiler and Lehman. Each atom starts with a number for its element,
scrambled by a fixed function that spreads nearby numbers far apart. In
each round it replaces its number by the scrambled combination of its own
number and the sum of its neighbours' scrambled numbers, so that after
$m$ rounds the number describes the atom's surroundings $m$ bonds deep,
and the molecule's number is the sum over its atoms. In a ring of six
every atom has two neighbours with two neighbours each, all round; at the
ends of a chain an atom has one, and its number differs from the first
round on. The scheme is not perfect. Graphs in which every atom has the
same number of neighbours of the same kinds, such as a triangular prism of
six atoms and a hexagon of six whose opposite atoms are also joined, look
alike to every round and get the same number, so two species can, rarely,
share a hash.
\end{toolbox}

\noindent The scrambling is the finishing step of the generator
splitmix64, a few multiplications and shifts of 64-bit whole numbers whose
sums are kept to their lowest 64 bits, dropping whatever passes
$2^{64}$; four rounds are used. \texttt{ReactionLog}
names each species by its formula, and a further isomer of a formula
already seen by the formula with a prime, A$_2$B$'$, in order of
appearance.

\subsection*{Checks}

The tracker's first frame, with no hysteresis yet, must find the pairs
closer than the sum of their radii, as ASE's \texttt{neighbor\_list} does
with its \texttt{natural\_cutoffs}. On 80 atoms of carbon and hydrogen at
random in a periodic cell, with $r_{\mathrm{on}}$ the sum of the radii,
the two lists are identical,
and both equal the list from testing every pair. Over 40 frames of atoms
moving at random, the tracker's bonds equal, frame by frame, those of a
direct transcription of the rule. On invented trajectories whose events
are known, it finds them all: of two molecules N$_2$, one pulled apart gives N$_2$
$\to$ N + N; one of the freed atoms brought to the other molecule gives N
+ N$_2$ $\to$ N$_3$; a ring of six carbon atoms, one of whose bonds is
stretched beyond $r_{\mathrm{off}}$, becomes a chain whose hash differs,
C$_6$ $\to$ C$_6'$. Each is a test of \texttt{mdlab}.

\begin{figure}[htb]
\centering
\includegraphics{ch16_observables/bonds}
\caption{The A-B gas: 250 molecules in a cell of side \SI{63}{\angstrom}
under Langevin friction of \SI{1}{\per\pico\second} for
\SI{100}{\pico\second}. (a) Bonds formed or broken per picosecond at
\SI{2000}{\kelvin}, counted with a single cutoff at $r_{\mathrm{on}}$ and
with the band from $r_{\mathrm{on}}$ to $r_{\mathrm{off}}$, against the
least lifetime an event must have to count. (b) The numbers of molecules
AB, of free A and of larger complexes at \SI{2000}{\kelvin} against time.
(c) The equilibrium constant $n_{\mathrm{AB}}V/n_{\mathrm A}n_{\mathrm B}$
over the last \SI{90}{\pico\second} with the band (circles) and with the
single cutoff (squares), with the standard error of nine blocks, against
$\int_0^{r_{\mathrm{on}}}4\pi r^2\mathrm{e}^{-\beta\varphi}\,\dd r$
(dashed).}
\label{fig:ob-bonds}
\end{figure}

\begin{keyidea}
A bond is assigned from the distance, with a band of hysteresis, forming
below $r_{\mathrm{on}}$ and breaking above $r_{\mathrm{off}}$, and events
are counted only for bonds and gaps that last a least time. Bonds are whole-number
keys, and those formed and broken are differences of sorted lists.
Molecules are the connected components of the graph of bonds, and a
species is named by its formula and a hash of its graph that does not
depend on the numbering.
\end{keyidea}

\notebook{16}{watch bonds flicker with one cutoff and settle with a band}

\begin{selfcheck}
In a run of 1000 atoms, the bond between atoms 17 and 342 has the key
17\,342. Which bond has the key 342\,017?
\end{selfcheck}
